\calculusName{PermutativeLambek Calculus}   
\calculusAcronym{\PLambekCalc}     
\calculusLogic{Substructural Logics} 
\calculusLogicOrder{Propositional}
\calculusType{Sequent Calculus}   
\calculusYear{1985}   
\calculusAuthor{Joachim Lambek}

\entryTitle{Permutative Lambek Calculus}        
\entryAuthor{Valeria de Paiva}  

%\etag{Non-Commutative}

\maketitle

\begin{entry}{PermLambekCalc}  

%\newcommand{\llimp}[0]{\leftharpoonup}
%\newcommand{\rlimp}[0]{\rightharpoonup}    
  
\begin{calculus}
\[
\begin{array}{ccccccccccc}
  & \quad & \infer[ax]{A \vdash A}{} & \quad & \infer[cut]{\Gamma_1,\Gamma_2,\Gamma_3 \vdash C}{\Gamma_1 \vdash A & \Gamma_2,A,\Gamma_3 \vdash C}
  \\[8pt]
  \infer[I_r]{\cdot \vdash I}{}
  & \quad &
  \infer[I_l]
        {\Gamma_1,I,\Gamma_2 \vdash A}
        {\Gamma_1,\Gamma_2 \vdash A}
  & \quad &  
  \infer[\otimes_r]
        {\Gamma_1,\Gamma_2 \vdash A \otimes B}
        {\Gamma_1 \vdash A & \Gamma_2 \vdash B}
  & \quad &
  \infer[\otimes_l]
        {\Gamma_1,A \otimes B,\Gamma_2 \vdash C}
        {\Gamma_1,A,B,\Gamma_2 \vdash C}
  \\[8pt]
  \infer[\limp_r]{\Gamma \vdash A \limp B}{\Gamma,A \vdash B}
  & \quad &
  \infer[\limp_l]{\Gamma_1,A \limp B,\Gamma_2,\Gamma_3 \vdash C}{\Gamma_1 \vdash A & \Gamma_2, B, \Gamma_3 \vdash C}
  & \quad 
\end{array}
\]
\end{calculus}

\begin{clarifications}
The Permutative Lambek Calculus  (LP) described here was introduced by Johan van Benthem to connect the original Lambek calculus to then new developments in Categorial Grammars.  Lambek  introduced two calculi, one in 1958, already in this Encyclopedia and one in 1961\cite{lambek1961}where tensor is not even associative. 
  \end{clarifications}

\begin{history}
The system now known as the basic Lambek Calculus was introduced in
1958 by Joachim Lambek as the ``Syntactic Calculus''~\cite{lambek1958}.
Lambek's motivation was to ``to obtain an effective rule (or
algorithm) for distinguishing sentences from non-sentences, which
works not only for the formal languages of interest to the
mathematical logician, but also for natural languages $[\ldots]$'', as
explained by Moortgat in~\cite{moortgat2010}.  After a long period of
neglect, around the middle 1980s the Syntactic Calculus, now called
the Lambek Calculus was taken up by logicians interested in
Computational Linguistics, especially van Benthem, Buszkowski and
Moortgat. They realized that a computational semantics for categorical
derivations along the lines of the Curry-Howard proofs-as-programs
interpretation would provide  a ``parsing-as-deduction''
paradigm and a powerful tool to study ``logical'' derivational
semantics. Around the same time, the introduction of Linear Logic~\iref{LL}, by
Jean-Yves Girard also gave a new impulse to the work in
Categorical Grammars, now called 
\end{history}

%% \begin{technicalities}
%% ...
%% \end{technicalities}

\end{entry}
